\section{Introduction}

Can chart images help predict stock returns beyond standard tabular representations of price and volume? Recent work has revived this classical technical-analysis question by converting historical OHLCV data into images and training convolutional neural networks on chart-like inputs \citep{jiang2023reimagining}. The intuitive appeal is clear: if investors visually recognize price patterns, an image model may extract shape-based signals that are difficult to hand-engineer.

This paper argues that the central question is broader. When CNNs and numerical machine-learning models are restricted to exactly the same 20-day OHLCV history, image representation does not deliver a stable single-model advantage over strong nonlinear numerical learners. In our China A-share sample, a chart-image CNN produces substantial return predictability, but HistGB, a stabilized bagged MLP, and especially a 1D-CNN trained on the same OHLCV sequence match or exceed its long-short performance. This finding weakens a simple ``image superiority'' story, but it opens a more interesting puzzle: why do different nonlinear learners trained on the same information set produce distinct, economically meaningful stock rankings?

We study this puzzle in six steps. First, we construct a matched-information benchmark in which all models use the same historical OHLCV inputs and are evaluated on the same out-of-sample panel. Second, we examine explicit chart indicators and show that they primarily improve linear models, suggesting that feature engineering changes the accessibility of information rather than the information set itself. Third, we document substantial complementarity across CNN, HistGB, and MLP predictions. Fourth, we introduce a temporal numerical benchmark: a 1D-CNN trained on the same 120D inputs reshaped into a $20\times 6$ sequence. Fifth, we decompose predictions into shared and model-specific components. Finally, we connect these components to recent price-path structures and to the horizon and state dependence of future returns.

The evidence points to a model-specific view of price-history alpha. Shared signals are strong and largely reflect recent return and volume consensus. In the three-model CNN--MLP-4--HistGB decomposition, CNN-specific and MLP-4-specific signals survive mutual controls. The new temporal benchmark sharpens this interpretation: temporal convolution absorbs a large share of the CNN-versus-flattened-model residual, but it does not fully span the chart CNN. In the separate pairwise CNN--1D-CNN residual tests, chart-CNN content conditional on 1D-CNN is more sensitive to near-end local path and volume structure and is realized chiefly at short horizons, whereas 1D-CNN-specific content conditional on CNN is more persistent and stronger in high-pressure states. Thus, the same price history can be parsed into distinct and complementary predictive signals by different nonlinear function classes and temporal encodings.

This paper contributes to three literatures. First, it contributes to machine learning for return prediction by separating information-set effects from model-class and temporal-encoding effects \citep{gu2020empirical}. Second, it contributes to the empirical literature on technical analysis by translating model-specific predictions into interpretable price-path structures. Third, it contributes to work on model combination and forecast diversity by showing that economically useful diversity can arise even when models share identical raw inputs. The current temporal benchmark is intentionally minimal: it establishes that local temporal convolution is a strong numerical alternative, while leaving richer TCN, recurrent, and transformer specifications as robustness extensions.
